\section{The cubic character field}\label{sec:cub}

The Frobenius--Perron dimension field means the field generated by
all simple Frobenius--Perron dimensions. Rank-four fusion rings are commutative by
Lemma~\ref{core:lem:rank-four-algebra}. Suppose a complex fusion categorification exists
and its dimension field $K$ has degree three. The ordinary rational algebra is then $\Q\times K$,
and character cyclotomicity makes $K$ cyclic and totally real
\cite[Corollary~8.53]{ENO}. Every character is consequently real.
The based involution acts by complex conjugation on characters, so it
is the identity: self-duality follows automatically.
Its Frobenius--Perron codegree $D$ is also
cubic. Indeed, if $D$ were rational, the reciprocal-codegree identity
would give $(f-1)D=3f$ for the rational codegree $f$, hence
$(f,D)=(2,6)$ or $(4,4)$. The corresponding order discriminants are
$432$ and $256$. The first is not a square, whereas a cyclic cubic
field has square discriminant; the second cannot be divisible by the
square of a nontrivial odd cubic conductor. Both are impossible.

Central induction places $D$ in the field generated by the signed
spherical dimensions. That field has degree at most four, so it has
degree three and equals $K$. Thus the signed rational algebra is also
$\Q\times K$, with its dimension character in the cubic factor.
The setup below introduces no additional pivotality or equality-of-fields
assumption.

\subsection{The cubic signed dimension field}

Let $\C$ be a self-dual complex fusion category of rank four,
and let $\widetilde{\C}$ be its sphericalization, with central
sign $\varepsilon$ of dimension $-1$; see \cite[Section~5.1]{Ostrik3} and
\cite[Section~5.3]{Bartlett}.  Its signed rational algebra is
\[
 T_{\Q}
 =\bigl(K_0(\widetilde{\C})/(\varepsilon+1)\bigr)
      \otimes\Q.
\]
Assume that $T_{\Q}\cong\Q\times K$, where $K$ is a
totally real cyclic cubic field, and that the spherical dimension
character belongs to the cubic factor.  Choose dual-compatible lifts
$X_i^+$ of the four simples, with $X_0^+=\mathbf1$, and write
$X_i^-=\varepsilon\otimes X_i^+$,
$x_i=\dim X_i^+$, and $E=\sum_i x_i^2$.  Let $\rho$ denote the
rational character of $T_{\Q}$ and let
$g=\sum_i\rho(X_i^+)^2$ be its codegree.  Both rank-four algebras are commutative by
Lemma~\ref{core:lem:rank-four-algebra}.  The sphericalization algebra, their direct sum, is
therefore commutative as well.  The Frobenius trace identity is
\begin{equation}\label{cub:eq:cubic-corrected-frobenius}
 \operatorname{Tr}_{K/\Q}(x_i/E)
       =\delta_{i0}-\frac{\rho(X_i^+)}g.
\end{equation}
To see this identity, let $\tau$ be the coefficient of the unit.
In the orthonormal signed basis, the regular trace of multiplication by
$a$ equals $\tau(a C_T)$, where $C_T$ is the Casimir.  The rational and
cubic components of $C_T$ are $g$ and $E$.  Evaluating this formula on
$X_i^+ C_T^{-1}$ gives the displayed identity.

Let $f_1,\ldots,f_4$ be the formal codegrees of $K_0(\C)$,
and let $g,E_1,E_2,E_3$ be those of the signed algebra.  Put
\[
 u=\sum_{i=1}^4f_i^{-2},\qquad
 v=g^{-2}+\sum_{i=1}^3 E_i^{-2}.
\]
As usual, these numbers are rational.  The dimension-field theorem
\cite[Corollary~2.15]{Ostrik3} places all the displayed codegrees in $K$.
Throughout, codegrees are counted with their character multiplicities.
The two reciprocal-codegree sums are one, and $f,g\ge2$ whenever the
ordinary rational character $f$ is introduced below.

\subsection{A positive central measure}
\begin{proposition}\label{cub:prop:cubic-corrected-average}
There is a positive, finitely supported measure $\xi$ on the unit circle,
of mass three, such that
\[
 \int z^n\,d\xi(z)=
 \begin{cases}1/2,&n\text{ odd},\\1,&n\text{ even},\end{cases}
 \qquad K\cap\Q(\zeta_n)=\Q.
\]
If the ordinary algebra is $\Q\times K$, with rational codegree
$f$, its known atoms satisfy
\begin{align}
 \xi(\{1\})&\ge\frac34(2u+v)+\frac{f-2}{f^2}
                          +\frac{g-2}{2g^2},\label{cub:eq:positive-atom}\\
 \xi(\{-1\})&\ge\frac34v+\frac{g-2}{2g^2}\ge\frac14.
 \label{cub:eq:negative-atom}
\end{align}
\end{proposition}
\begin{proof}
Let $F:Z(\widetilde{\C})\to\widetilde{\C}$ be forgetful
and $I$ its right adjoint. Use the canonical normalized center matrix
$S=s/(2E)$. The rational character $\rho$ corresponds to a self-dual
constituent $B\ne1$ of $I(1)$ by
\cite[Example~2.9 and Theorem~2.13]{Ostrik3}. Put
\[
 a_A=S_{0A}=\frac{\dim A}{2E}\in K,\qquad
 b_A=S_{BA}=\frac{\rho(F(A))}{2g}\in\Q.
\]
Orthogonality gives
\begin{equation}\label{cub:eq:cubic-corrected-row-norms}
 \sum_Aa_A^2=\sum_Ab_A^2=1,\qquad \sum_Aa_Ab_A=0.
\end{equation}
In a common cyclotomic Galois field $\Omega$, write
$G=\operatorname{Gal}(\Omega/\Q)$. Its signed permutations
$\widehat\sigma$ satisfy
\begin{equation}\label{cub:eq:cubic-corrected-galois}
 \sigma(a_A)=\epsilon_\sigma(A)a_{\widehat\sigma A},\quad
 b_A=\epsilon_\sigma(A)b_{\widehat\sigma A},\quad
 \theta_{\widehat\sigma A}=\sigma^2(\theta_A).
\end{equation}
The last identity has no scalar anomaly for the canonical center
normalization \cite[Lemma~4.2 and Section~4.1]{DLN}. Define
\[
 \xi=\frac1{|G|}\sum_{\sigma\in G}\sum_A
 \bigl(3\sigma^2(a_A)^2+\sigma^2(a_A)b_A\bigr)
 \operatorname{Dirac}_{\sigma^2(\theta_A)}.
\]
Here $\sigma^2$ means composition. Equation
\eqref{cub:eq:cubic-corrected-row-norms} gives mass three.

For an allowed $n$, the restrictions of $G$ to $K$ and
$\Q(\zeta_n)$ are independent. Squaring permutes
$\operatorname{Gal}(K/\Q)\cong C_3$; no corresponding
surjectivity is required on the indicator field.
Apply \eqref{core:central-indicator} and
\eqref{cub:eq:cubic-corrected-frobenius}. The moment of the first
term of $\xi$ is
\[
 \frac{1-\nu_n(\varepsilon)}2
 -\frac1{2g}\operatorname{Av}_{\sigma\in G}
 \sigma^2\!\left(\sum_i\rho(X_i^+)
          [\nu_n(X_i^+)-\nu_n(X_i^-)]\right).
\]
The sum in parentheses is $2g\sum_Aa_Ab_A\theta_A^n$ by the same
center identity. The second term of $\xi$ cancels this residual.
The sign object has indicator $-1$ for even $n$ and zero for odd $n$:
its even invariant space is one-dimensional and rotation includes its
inverse pivotal map $-1$, whereas the odd invariant space vanishes.
This proves the moment formula.

Reindexing each summand by $W=\widehat\sigma A$ fixes the twist
$\theta_W$ and gives its individual weight
\begin{equation}\label{cub:eq:cubic-corrected-weight}
 w_W=\operatorname{Tr}_{K/\Q}(a_W^2)
              +\frac{b_W}{3}\operatorname{Tr}_{K/\Q}(a_W).
\end{equation}
In particular the measure does not depend on enlarging $\Omega$.
Let $h_W=[F(W):1]-[F(W):\varepsilon]$. The Frobenius trace identity
says $\operatorname{Tr}(a_W)=h_W/2-b_W$.
Outside $I(1)\cup I(\varepsilon)$ one has $h_W=0$, so
$w_W=\operatorname{Tr}(a_W^2)-b_W^2/3\ge0$.
For $W\subset I(1)$ put $t_W=\operatorname{Tr}(a_W)>0$.
All conjugate codegrees are positive, and $b_W=1/2-t_W$; hence
\[
 w_W\ge\frac{t_W^2}{3}+\frac{(1/2-t_W)t_W}{3}
       =\frac{t_W}{6}>0.
\]
Its sign translate has the same weight. Thus positivity holds
pointwise, without grouping atoms.

For the four ordinary constituents the $b$-values are
$(f-3)/(2f),1/(2f),1/(2f),1/(2f)$; the signed values replace $f$
by $g$. This follows from $a_A=1/(2f_\chi)$ and
$b_A=1/2-\operatorname{Tr}(a_A)$. Summing their corrections in
\eqref{cub:eq:cubic-corrected-weight} gives $(f-2)/(2f^2)$ and
$(g-2)/(2g^2)$, respectively. By
Lemma~\ref{core:lem:translate-twists}, the
ordinary sign translates have twist $+1$ and the signed translates
have twist $-1$. Summing their weights gives
\eqref{cub:eq:positive-atom}--\eqref{cub:eq:negative-atom}; the final
bound uses $v\ge g^{-2}+(1-g^{-1})^2/3$.
\end{proof}

The pushforward under $z\mapsto z^2$ is the resulting squared-twist
measure. If $K\cap\Q(\zeta_{2j})=\Q$ for $1\le j\le m$,
its mass is three, its first $m$ moments are one, and its atom at one
has weight at least $3(u+v)/2$. The Fej\'er bound in
Lemma~\ref{core:lem:measure-tools} therefore gives
\begin{equation}\label{cub:eq:cubic-corrected-six-sevenths}
 u+v\le\frac{2(m+3)}{3(m+1)},\qquad u+v\le\frac67\quad(m=6).
\end{equation}
Likewise $\xi-\tfrac14\operatorname{Dirac}_{-1}$ has mass $11/4$
and all available moments equal to $3/4$. Thus
\begin{equation}\label{cub:eq:unsquared-fejer}
 \frac34(2u+v)+\frac{f-2}{f^2}+\frac{g-2}{2g^2}
 \le\frac34+\frac2{m+1}
\end{equation}
when the first $m$ indicator fields are disjoint from $K$.

\subsection{Periodic moments}
\begin{proposition}\label{cub:prop:cubic-periodic}
Suppose the ordinary and signed rational algebras are $\Q\times K$
as above. If $F$ is the conductor of $K$, then
\begin{equation}\label{cub:eq:cubic-periodic}
 \mathcal L:=\frac32(u+v)+\frac{f-2}{f^2}+\frac{g-2}{g^2}
          \le1+\frac2F\le\frac97.
\end{equation}
\end{proposition}
\begin{proof}
A nontrivial cyclic cubic conductor is odd and at least seven.
Since $[K:\Q]=3$, the field $K$ is disjoint from
$\Q(\zeta_n)$ whenever $F\nmid n$.
Restrict $\xi$ to the $2F$th roots and call the restriction $\xi_0$.
For $F\nmid r$, average its moments with exponents $2Fj+r$ over
$0\le j<M$ and let $M\to\infty$. The geometric average is zero off
the $2F$th roots and one on them, so $\xi_0$ has the same prescribed
$r$th moment. Its mass is at most three.
Subtract $\tfrac34\operatorname{Dirac}_1+
\tfrac14\operatorname{Dirac}_{-1}$. The residual Fourier transform
on this cyclic group is supported on $\{0,F\}$, so its weights are
constant on each parity class; call them $a,b$.
Each class has roots other than $\pm1$, at which these are weights
of the positive measure $\xi_0$. Hence $a,b\ge0$ and
\[
 1+F(a+b)\le3,\qquad
 \xi(\{1\})+\xi(\{-1\})=1+a+b\le1+\frac2F.
\]
Now use \eqref{cub:eq:positive-atom}--\eqref{cub:eq:negative-atom}.
\end{proof}

\subsection{The rational codegrees and cubic norms}
The two rational codegrees $f,g$ are integers at least two. The
ordinary value $f=2$ is excluded by
Proposition~\ref{inv:prop:cubic-rational-two-impossible}.

\begin{lemma}\label{cub:lem:cubic-square-norm-schur}
For either algebra, let $r$ be its rational codegree and $H$ a codegree
from its cubic character orbit. There are positive integers $A,s$ such
that
\[
 \operatorname{Tr}(H)=rA,\qquad
 \operatorname{Norm}(H)=r^3s^2,
\]
and $H/r$ has minimal polynomial
\begin{equation}\label{cub:eq:cubic-normalized-polynomial}
 z^3-Az^2+(r-1)s^2z-s^2.
\end{equation}
If $r\ge3$, the inverse-square sum of the four codegrees satisfies
\begin{equation}\label{cub:eq:cubic-uniform-schur}
 r^{-2}+\sum_{\sigma:K\hookrightarrow\R}\sigma(H)^{-2}
 \ge r^{-2}+\tfrac12(1-r^{-1})^2.
\end{equation}
\end{lemma}

\begin{proof}
First $H$ has degree three. If it were rational, the reciprocal-codegree
identity would give $(r-1)H=3r$, so $(r,H)=(2,6)$ or $(4,4)$.
The order discriminant is the product of its four formal codegrees.
The first case gives $432$, which is not a square, contrary to the
cyclic cubic field discriminant. The second gives $256$; its square
root would have to be divisible by the conductor $F_K$, whereas a
nontrivial cyclic cubic conductor is odd and cannot divide sixteen.

Write the minimal polynomial of $H$ as
\[
 z^3-a z^2+(r-1)h z-rh.
\]
The reciprocal identity gives this form with positive integers $a,h$.
The d-number criterion gives $rh\mid a^3$ and
$(rh)^2\mid((r-1)h)^3$. Thus $r^2\mid h$ and $r\mid a$.
The order discriminant $r\operatorname{Norm}(H)=r^2h$ is a square,
so $h=r^2s^2$ and $a=rA$ for positive integers $A,s$. This proves
\eqref{cub:eq:cubic-normalized-polynomial}. For the signed algebra the
same argument applies: $2H$ is an actual formal codegree of the
sphericalization, and the d-number property is unchanged by nonzero
rational scaling.

Put $k=r-1$ and $n=s^2$. The three positive reciprocal roots of
\eqref{cub:eq:cubic-normalized-polynomial} have sum $k$, pairwise-product
sum $A/n$, and product $1/n$. Schur's inequality yields
\[
 4kA\le k^3n+9.
\]
If $k\ge3$ and $4A>k^2n$, the positive integer $4A-k^2n$ is at least
four when $k$ or $s$ is even, and at least three otherwise. This
contradicts the displayed inequality except possibly at $k=3$ with
equality. Equality in Schur for positive numbers makes all three roots
equal, contradicting their cubic degree. Hence $A\le k^2n/4$.
For $k=2$, Schur gives $A\le n+1$. At $A=n+1$ the discriminant is
$n(n-4)=s^2(s^2-4)$, which cannot be a nonzero square. Thus $A\le n$
in this case too. Finally the inverse-square sum is
\[
 \frac{1+k^2-2A/n}{r^2}
 \ge\frac1{r^2}+\frac{k^2}{2r^2},
\]
as claimed.
\end{proof}

\begin{theorem}[Direct cubic reduction]\label{cub:thm:cubic-reduction}
The only possibilities are
\[
 (F,f,g)=(7,3,3),\quad(7,4,4),\quad(9,3,3),
 \qquad \operatorname{Norm}(E)=\operatorname{Norm}(D).
\]
In the conductor-nine case $u=v=1/3$, the ordinary and signed cubic
codegrees agree as multisets, and the category is spherical.
\end{theorem}
\begin{proof}
For a constituent $A$ from the ordinary cubic orbit,
$a_A=1/(2D_i)$ gives $b_A=1/(2f)$ in this measure.
Since $b_A=\rho(F(A))/(2g)$ and $\rho$ is integral-valued,
$g/f\in\ZZ_{>0}$.
Put $a=\operatorname{Tr}(D)$, $b=\operatorname{Tr}(E)$,
$c=\operatorname{Norm}(D)$ and
$k=\operatorname{Norm}(E)/c\in\ZZ_{>0}$, using central
integrality. Casimir trace domination gives $b+g\le a+f$.
For
\[
 H(r)=\frac12-\frac1r+\frac{3}{2r^2},\qquad
 J(r)=1-\frac2r+\frac2{r^2},
\]
Lemma~\ref{cub:lem:cubic-square-norm-schur} and the reciprocal
identities give
\[
 u\ge H(f),\quad u=J(f)-2a/c,\quad v=J(g)-2b/(kc),
 \quad v-u\ge J(g)-J(f).
\]
If $g\ge2f$, substitution in \eqref{cub:eq:cubic-periodic} yields
\[
 \mathcal L\ge\frac32+\frac1f-\frac1{2f^2}-\frac2g+\frac1{g^2}
 \ge\frac32-\frac1{4f^2}\ge\frac{53}{36}>\frac97,
\]
a contradiction. Therefore $g=f$. Now
\[
 \mathcal L\ge\frac32-\frac1f+\frac1{2f^2}.
\]
This increases for $f>1$ and equals $33/25>9/7$ at $f=5$.
Consequently $f=g\in\{3,4\}$.

Lemma~\ref{cub:lem:cubic-square-norm-schur} makes $k$ an integer
square. If $k>1$, then $k\ge4$, and $b\le a$ implies
$v\ge\tfrac34J(f)+\tfrac14u$. Thus
\[
 \mathcal L\ge\frac{33}{16}-\frac{17}{8f}+\frac{17}{16f^2}
 \ge\frac{53}{36}>\frac97,
\]
again impossible. Hence $k=1$.
At $f=3$ the lower bound for $\mathcal L$ is $11/9$, forcing
$F\le9$; at $f=4$ it is $41/32$, forcing $F\le64/9$.
The only cyclic cubic conductors in these ranges are seven and nine,
as follows from primitive Dirichlet characters of order three.
This gives the stated list.

At $(F,f)=(9,3)$ equality holds throughout, so $u=v=1/3$.
Equality of norms and inverse-square sums gives
$\operatorname{Tr}(E)=\operatorname{Tr}(D)$, and the reciprocal
identity identifies the remaining coefficient of the two cubic
polynomials. Their roots therefore agree. Choose a signed character
with codegree $D$ and apply spectral saturation in
Lemma~\ref{core:lem:family-induction} to obtain a spherical section.
\end{proof}

\subsection{The conductor-nine ring}

\begin{theorem}\label{cub:thm:cubic-nine-ring}
Let $\C$ be a self-dual rank-four complex fusion category whose
Frobenius--Perron codegree field has conductor nine. Its Grothendieck
ring, up to relabeling, has basis $\mathbf1,X,Y,Z$ and products
\begin{equation}\label{cub:eq:cubic-nine-ring}
\begin{aligned}
 X^2&=\mathbf1+Z,& XY&=Y+Z,& XZ&=X+Y,\\
 Y^2&=\mathbf1+X+Y+Z,& YZ&=X+Y+Z,& Z^2&=\mathbf1+Y+Z.
\end{aligned}
\end{equation}
This ring occurs as the even part of $\mathrm{SU}(2)_7$, and every
complex fusion category with this Grothendieck ring admits a spherical
structure.
\end{theorem}

\begin{proof}
Theorem~\ref{cub:thm:cubic-reduction} gives $f=3$ and $u=1/3$.
Write $n=s^2$ in Lemma~\ref{cub:lem:cubic-square-norm-schur}. The inverse-square
sum is $(5-2A/n)/9$, so $A=n$. Consequently $D/3$ has polynomial
\[
 z^3-nz^2+2nz-n.
\]
Its discriminant is $n^2(4n-27)$. The cubic field is cyclic, so this
is a nonzero square. Thus $r^2=4s^2-27$ for a positive integer $r$,
and
\[
 (2s-r)(2s+r)=27.
\]
The positive factor pairs $(1,27)$ and $(3,9)$ give $s=7$ and $s=3$.
At $s=7$ the polynomial discriminant is $7^4\cdot13^2$, whereas
$K$ has discriminant $9^2$, which must divide the discriminant of any
integral power basis. This excludes $s=7$. Hence $s=3$, $n=A=9$,
$\Tr(D)=27$, and
\begin{equation}\label{cub:eq:cubic-nine-codegree}
 D^3-27D^2+162D-243=0.
\end{equation}

The rational character has codegree three, so its nonunit values are
two signs and a zero. Orthogonality to the positive dimension row
rules out both signs positive; two negative signs would give two
positive simple dimensions summing to one. After relabeling its row
is $(1,1,-1,0)$. Frobenius reciprocity makes the three-index structure
tensor symmetric. Applying this character to each product therefore
puts the table in the form
\begin{equation}\label{cub:eq:cubic-three-normal-form}
\begin{aligned}
 X^2&=1+aX+aY+cZ,& XY&=aX+(a+1)Y+cZ,\\
 XZ&=cX+cY+pZ,&Y^2&=1+(a+1)X+(a+1)Y+cZ,\\
 YZ&=cX+cY+(p+1)Z,&Z^2&=1+pX+(p+1)Y+jZ,
\end{aligned}
\end{equation}
where $a,c,p,j$ are nonnegative integers. One associativity equation,
obtained by comparing the $X$ coefficients of $X(Z^2)$ and $(XZ)Z$,
is
\begin{equation}\label{cub:eq:cubic-three-assoc}
 cj=2c^2+p^2-a(2p+1)-1.
\end{equation}
The trace of the Casimir is the sum of the squares of all structure
coefficients, and also the sum of the four codegrees. Substituting
\eqref{cub:eq:cubic-three-normal-form} and using $3+\Tr(D)=30$ gives
\[
 17+8a(a+1)+6p(p+1)+12c^2+j^2=30.
\]
Thus $a=0$, and the remaining equality has only the possibilities
$(c,p,j)=(1,0,1)$ and $(0,1,1)$. Indeed, $c,p\le1$; when both vanish
$j^2=13$, and when both equal one the left side already exceeds thirty.
Both surviving choices satisfy~\eqref{cub:eq:cubic-three-assoc}.
In the second table $X^2=1$, $XY=Y$, and $XZ=Z$, so there is a second
rational character $(1,-1,0,0)$. This contradicts the assumed rational
algebra $\Q\times K$. The first choice is precisely
\eqref{cub:eq:cubic-nine-ring}.

For existence use the even part of $\mathrm{SU}(2)_7$
\cite[Section~5.3.10]{RSW}. Its simple labels are $V_0,V_2,V_4,V_6$.
The truncated tensor-product rule is
\[
 V_i V_j=\sum_{\substack{|i-j|\le\ell\le\min(i+j,14-i-j)\\
                           \ell\equiv i+j\ (\mathrm{mod}\ 2)}}V_\ell.
\]
Taking $X=V_6$, $Y=V_4$, and $Z=V_2$ yields the displayed products.
Their Frobenius--Perron dimensions lie in $\Q(\zeta_9+\zeta_9^{-1})$,
and~\eqref{cub:eq:cubic-nine-codegree} has degree three, so this is exactly
the codegree field. The sphericality assertion follows from
Theorem~\ref{cub:thm:cubic-reduction}.
\end{proof}
\subsection{An inert-prime obstruction for a rational codegree}

\begin{lemma}\label{cub:lem:rational-prime-inert}
Let $R$ be a commutative self-dual based ring with
$R_{\Q}\cong\Q\times K$, and suppose its rational
character has formal codegree equal to a prime $p$.
Then $K\otimes\Q_p$ cannot be an unramified field extension
of $\Q_p$ of degree greater than one.
\end{lemma}
\begin{proof}
Suppose that $L=K\otimes\Q_p$ is such an extension.
Put $A=R\otimes\ZZ_p$ and let $B$ be its image in $L$.
Thus $B$ is a finite $\ZZ_p$-order in the valuation ring
$\mathcal O_L$. Write $\rho_i$ for the integral values of the
rational character, with $\rho_0=1$. Its primitive idempotent is

\[
 e=\frac1p\sum_i\rho_i b_i=(1,0)
 \quad\text{in }\Q_p\times L.
\]

Consequently $A\cap(\Q_p\times0)=p\ZZ_p\times0$:
the coefficient of $b_0$ in $te$ is $t/p$, and all other
coefficients are integral once $t\in p\ZZ_p$.
It follows that the first coordinate modulo $p$ defines a
well-defined surjective ring map $B\to\mathbb F_p$, and

\[
 A=\ZZ_p\times_{\mathbb F_p}B.
\]

Indeed, if $(a_0,b)\in A$ and $a\equiv a_0\pmod p$, then
$(a,b)=(a_0,b)+(a-a_0,0)\in A$.
The order $B$ is local: its integral normalization $\mathcal O_L$
has a unique maximal ideal, and lying-over gives the same assertion
for $B$. Its displayed quotient is therefore its residue field.

The unit-coefficient pairing on $A$ is perfect and associative.
It gives an $A$-module isomorphism
$A\cong\operatorname{Hom}_{\ZZ_p}(A,\ZZ_p)$. As $A$ is finite free over
the discrete valuation ring $\ZZ_p$, this identifies its
canonical module with $A$, so $A$ is Gorenstein.

The fiber-product theorem~\cite[Fact~2.6 and Corollary~2.7]{NTV}
states that a
nontrivial fiber product of noetherian local rings over their common
residue field can be Gorenstein only if both factors are regular
local rings of dimension one. Thus $B$ is a discrete valuation ring.
It is then integrally closed in $L$, and hence $B=\mathcal O_L$.
But the latter has residue field $\mathbb F_{p^{[L:\Q_p]}}$,
contradicting the residue field $\mathbb F_p$ of $B$.
\end{proof}

\begin{corollary}\label{cub:cor:cubic-seven-three-impossible}
There is no commutative self-dual based ring whose rational algebra
is $\Q\times\Q(\zeta_7+\zeta_7^{-1})$ and whose
rational formal codegree is three.
\end{corollary}
\begin{proof}
The polynomial $x^3+x^2-2x-1$ of
$\zeta_7+\zeta_7^{-1}$ has discriminant $49$ and has no root
modulo three: its reduction is $x^3+x^2+x-1$, taking nonzero
values at all three elements of $\mathbb F_3$. Thus three is
unramified and inert in this cubic field. Apply
Lemma~\ref{cub:lem:rational-prime-inert} with $p=3$.
\end{proof}


\subsection{Sphericality at rational codegree four}

\begin{lemma}\label{cub:lem:cubic-four-trace-gap}
In the remaining conductor-seven case, put
$a=\Tr(D)$, $b=\Tr(E)$, and $c=\operatorname{Norm}(D)
=\operatorname{Norm}(E)$. Then
\[
 a-b\in64\ZZ_{\ge0}.
\]
\end{lemma}
\begin{proof}
Let $B$ be the constituent of $I(\mathbf1)$ corresponding to the
rational signed character. Let $C_R$ and $C_T$ be the ordinary and
signed Casimir elements. We show that the ordinary folded and signed
forgetful coefficient vectors of $B$ are respectively
\[
 w=C_R/4,\qquad r=C_T/4.
\]
The signed cubic component of $F(B)$ is its dimension $E/4$.
Its unit coefficient is one, and its sign-object coefficient is zero,
because $B$ belongs to $I(\mathbf1)$ and the supports of $I(\mathbf1)$
and $I(\varepsilon)$ are disjoint. The Frobenius trace formula then
makes its signed rational component one, so its signed class is $C_T/4$.
For an ordinary cubic induction constituent $A_i$, the rational signed
row, already computed in Theorem~\ref{cub:thm:cubic-reduction},
has $S_{B,A_i}=1/8$. Symmetry and the modular character formula give
\[
 \chi_i(F(B))=\frac{S_{A_i,B}}{S_{A_i,0}}
             =\frac{1/8}{1/(2D_i)}=D_i/4.
\]
Its ordinary folded unit coefficient is likewise one, so the ordinary
rational component is one and its folded class is $C_R/4$.

These are two coefficient vectors of the same actual object. Thus
$w_i\ge|r_i|$ and $w_i\equiv r_i\pmod2$ for every $i$. In the
orthonormal based trace, the squared coefficient norm of a Casimir is
the sum of its formal codegrees. Consequently
\[
 a-b=16\bigl(\|w\|^2-\|r\|^2\bigr)\in64\ZZ_{\ge0}.
\]
\end{proof}

\begin{theorem}\label{cub:thm:cubic-four-spherical}
Every category in the conductor-seven, rational-codegree-four case
admits a spherical structure. Its signed and ordinary cubic codegree
polynomials coincide.
\end{theorem}
\begin{proof}
Retain $a,b,c$ from Lemma~\ref{cub:lem:cubic-four-trace-gap}. Here
\[
 u=\frac58-\frac{2a}{c},\qquad
 v=\frac58-\frac{2b}{c},\qquad
 u,v\ge\frac{11}{32}.
\]
The periodic inequality is $u+v\le29/42$. It follows that
\begin{equation}\label{cub:eq:cubic-four-trace-interval}
 \frac{187c}{1344}\le b\le a\le\frac{9c}{64},
 \qquad a-b\le\frac{c}{672}.
\end{equation}
If $a>b$, the trace-gap lemma gives $c\ge43008$.

For either triple of cubic codegrees let $H$ be its largest element
and put $x=1/H$. Since $H$ is at least a third of the trace,
\eqref{cub:eq:cubic-four-trace-interval} gives $0<x<1/1900$.
Let $q_1,q_2$ be the other two reciprocal codegrees. Their sum is
$3/4-x$. Also $u,v\le29/42-11/32=233/672$, so
\[
 q_1^2+q_2^2+x^2\le\frac{191}{672},\qquad
 (q_1-q_2)^2\le\frac1{168}+\frac32x-3x^2.
\]
The last upper bound is less than $(1/12-x)^2$, since the difference is
\[
 \frac1{1008}-\frac53x+4x^2
 >\frac1{1008}-\frac1{1140}>0.
\]
Therefore $1/3<q_1,q_2<5/12$. The two smaller codegrees of each
triple lie in $(12/5,3)$.

Align the largest signed codegree $E_{\max}$ with the ordinary
Frobenius--Perron codegree $D_{\rm FP}$. Central induction makes
$z=E_{\max}/D_{\rm FP}$ an algebraic integer: it gives integrality
against every conjugate of $D$, and one may then apply a field
embedding sending $E$ to $E_{\max}$. Equality of norms makes $z$
a totally positive unit of norm one. Spectral domination gives
$z\le1$ at the chosen embedding. Its other two conjugates are ratios
of numbers in $(12/5,3)$ and are less than $5/4$. Hence
\[
 3\le\Tr(z)<1+2\cdot\frac54=\frac72.
\]
The trace is an integer, so it is three; equality in the
arithmetic--geometric mean inequality forces every conjugate to be
one. Thus $z=1$, contradicting $a>b$.
We conclude $a=b$. Equality of norms and the reciprocal identity
then identifies the two cubic polynomials. Spectral saturation in
Lemma~\ref{core:lem:family-induction} supplies the spherical structure.
\end{proof}

\subsection{Excluding conductor seven}

The spherical section obtained above need not have positive dimensions
at its original embedding. Its construction does give an embedding
which sends every simple dimension to its Frobenius--Perron dimension.
Apply that embedding to the central indicator identities. Every central
dimension becomes positive, because its forgetful class has nonnegative
multiplicities, and every twist remains a root of unity. The first and
second indicators, respectively zero and $\pm1$ for a nonunit self-dual
simple, are unchanged. Thus the following argument uses the transformed
identities and does not assume that the original category is
pseudo-unitary.



\begin{theorem}\label{cub:thm:cubic-seven-excluded}
The conductor-seven, rational-codegree-four case cannot occur.
\end{theorem}
\begin{proof}
Since $u\ge11/32$ and $S\le\sqrt{3(D-1)}$, Lemma~\ref{core:lem:family-pivotal-fejer}, in the form
\eqref{core:three-term-unsimplified}, gives
\[
 \frac3{32}D\le6+2\sqrt{3(D-1)}.
\]
This forces $D<1500$. Indeed, for $D\ge1500$ its divided right side
is at most $1/250+1/(5\sqrt5)<3/32$.

Write $a=\Tr(D)$ and $n=s^2$ in
Lemma~\ref{cub:lem:cubic-square-norm-schur}, so that $a=4A$ and
$\operatorname{Norm}(D)=64n$. Since the codegree triples now agree,
$u=v$, and the periodic inequality gives $u\le29/84$. Together with
$u\ge11/32$ and $u=(10-2A/n)/16$, this yields
\begin{equation}\label{cub:eq:cubic-four-A-interval}
 \frac{47n}{21}\le A\le\frac{9n}{4}.
\end{equation}
Moreover the nonnegative integral vector $w=C_R/4$ from
Lemma~\ref{cub:lem:cubic-four-trace-gap} has unit coefficient one. Hence
\[
 a+4=16\|w\|^2,
 \qquad A=3+4\sum_{i\ne0}w_i^2\equiv3\pmod4.
\]

We sharpen the trace bound to obtain $s\le12$. If $D<100$, all three
ordinary cubic codegrees are at most $D$, so $a<300$ and
$n\le21a/188<34$. Otherwise put $x=1/D\le1/100$ and let $q_1,q_2$
be the reciprocals of the other two cubic codegrees. They satisfy
\[
 q_1+q_2=\frac34-x,
 \qquad (q_1-q_2)^2\le\frac1{336}+\frac32x-3x^2,
\]
using $u\le29/84$. The last upper bound is strictly less than
$(3/20-x)^2$, because their difference is
\[
 \frac9{400}-\frac1{336}-\frac95x+4x^2
 \ge\frac9{400}-\frac1{336}-\frac9{500}>0.
\]
Thus $q_1,q_2>3/10$; the other two codegrees are less than $10/3$.
Consequently
\[
 a<D+\frac{20}{3}<\frac{4520}{3},
 \qquad n\le\frac{21a}{188}<\frac{7910}{47}<169.
\]
In either case $s\le12$.

Finally the order discriminant is the product of the formal codegrees,
\[
 \operatorname{disc}(R)=4\operatorname{Norm}(D)=256s^2.
\]
The maximal order of $\Q\times K$ has discriminant $49$. Its index
over $R$ is therefore $16s/7$, an integer. Thus $7\mid s$, and
$1\le s\le12$ forces $s=7$, $n=49$. Inequality
\eqref{cub:eq:cubic-four-A-interval} now reads
$109\tfrac23\le A\le110\tfrac14$, so $A=110$. This contradicts
$A\equiv3\pmod4$.
\end{proof}

\begin{theorem}\label{cub:thm:cubic-classification}
A rank-four fusion ring whose Frobenius--Perron dimension
field has degree three admits a complex fusion categorification if and only
if it is the ring~\eqref{cub:eq:cubic-nine-ring}, the even part of
$\mathrm{SU}(2)_7$. Every complex fusion categorification of this ring
admits a spherical structure.
\end{theorem}
\begin{proof}
The periodic reduction and the inert-prime obstruction leave only the
conductor-nine case of Theorem~\ref{cub:thm:cubic-nine-ring} and the case
excluded in Theorem~\ref{cub:thm:cubic-seven-excluded}.
\end{proof}



